% Source: sources/returns/S6_04_MAPS_MANIFOLDS/sections/manifolds.tex:451–484
\paragraph{Theorem D.1.8 (Proved: the integral Pontryagin class).}
The pullback
\[
 \pi^*:H^4(X_p;\mathbb Z)\longrightarrow
 H^4(\widetilde X_p;\mathbb Z)
\]
is zero. In particular
$c_2(\widetilde X_p)=p_1(\widetilde X_p)=0$ integrally.
\begin{proof}
We apply the external-curve calculation in D.1.6 first downstairs.
The two mixed regular-fibre tori specialize to one external curve
and a signed sum of the two; their values under the primitive
functional $\Phi=(4,2,3,2)$ are $0$ and $2$.
The first external curve therefore has zero class in
$H_2(X_p)=\mathbb Z/d$, and the other, which we call $C$, has
class $\pm2$ and generates because $d$ is odd.
The covering has exactly $d$ lifts $C_j$ of $C$, each mapped
orientation-preservingly to $C$, and the transpose relations of
$B(t)$ identify all $[C_j]$ in $H_2(\widetilde W_p)$.
Their image is torsion of order $d$ in $H_2(\widetilde X_p)$
by the same mixed-torus calculation. Homology transfer consequently
takes the generator $[C]$ to
\[
 \operatorname{tr}[C]=\sum_j[C_j]=d[C_0]=0.
\]
For a finite orientation-preserving cover, Poincare duality
intertwines cohomological pullback with homological transfer;
this follows by lifting a transverse cycle or from the cap-product
chain identity with the lifted fundamental class. Thus $\pi^*$
vanishes in degree four. Naturality of Chern and Pontryagin
classes, applied to $T\widetilde X_p\simeq\pi^*TX_p$, proves
both assertions. For $d=1$ the target degree-four group already
vanishes.
\end{proof}
